El auge exponencial de la Inteligencia Artificial Generativa y los Modelos de Lenguaje de Gran Tamaño (\textit{Large Language Models}, LLMs) ha impulsado su rápida integración en el ámbito educativo como asistentes conversacionales y tutores virtuales. No obstante, su naturaleza intrínsecamente probabilística y no determinista introduce riesgos críticos de desinformación (alucinaciones factuales, sobreconfianza sintomática), amenazas a la integridad académica (plagio y resolución fraudulenta de evaluaciones) y deficiencias pedagógicas (solucionismo pasivo frente a andamiaje cognitivo).

En el ámbito de la Ingeniería del Software, evaluar estos sistemas plantea de forma directa el \textit{Problema del Oráculo} (\textit{The Oracle Problem}), dado que las respuestas en lenguaje natural no pueden validarse mediante aserciones booleanas deterministas clásicas.

Este Trabajo de Fin de Grado presenta una propuesta metodológica formal, sistemática y reproducible para el testing de calidad de chatbots educativos basados en LLMs. La metodología se fundamenta en los atributos de calidad del estándar internacional ISO/IEC 25010:2023, la Teoría del Andamiaje y Zona de Desarrollo Próximo (ZDP) de Vygotsky, y los modelos de feedback formativo de Hattie y Timperley. 

El marco define siete dimensiones de evaluación ortogonales (Corrección Factual, Control de Alucinaciones, Claridad Didáctica, Utilidad Pedagógica y Feedback, Robustez y Seguridad, Adaptación al Nivel y Seguimiento de Instrucciones), evaluadas mediante una matriz de rúbricas analíticas en escala forzada ($0-3$) para mitigar el sesgo de tendencia central. Asimismo, se formulan métricas cuantitativas agregadas y un Índice Global de Calidad Educativa ($IQE$) gobernado por el principio \textit{Safety-First}.

La metodología ha sido validada experimentalmente mediante la ejecución de una batería estructurada de 42 casos de prueba (126 ensayos) sobre tres perfiles conversacionales contrastados (\textit{Asistente Base}, \textit{Tutor Directo} y \textit{Tutor Socrático}), demostrando empíricamente que la ausencia de alineamiento pedagógico genera una tasa crítica de alucinaciones del $33.3\%$ y fallos de seguridad del $9.5\%$, mientras que el andamiaje socrático permite alcanzar una idoneidad formativa del $100\%$.

\palabrasclave{Modelos de Lenguaje de Gran Tamaño, Testing de Software, Chatbots Educativos, Calidad de Software, Alucinaciones, Andamiaje Socrático, ISO/IEC 25010, Inteligencia Artificial en Educación}
